\documentclass[aspectratio=169]{beamer}
\usetheme{Madrid}
\usecolortheme{default}
\usepackage[utf8]{inputenc}
\usepackage[spanish]{babel}

\title[Ruido en Perfiles Aerodinámicos]{Comparación experimental de algoritmos: \ Predicción de ruido en perfiles aerodinámicos}
\author[Ávila, González, Raya, Toledo]{Adrián G. Ávila Chávez \and Gustavo González De la Cruz \and Caleb de Jesús Raya Mendoza \and Mario D. Toledo González}
\institute[UGto]{Universidad de Guanajuato \ División de Ingenierías}
\date{\today}

\begin{document}

\begin{frame}
\titlepage
\end{frame}

\section{Metodología y Datos}

\begin{frame}{Análisis Exploratorio de Datos (EDA)}
\begin{itemize}
\item \textbf{Conjunto de datos:} 1,503 observaciones numéricas, obtenidas en túnel de viento anecoico (Perfil NACA 0012).
\item \textbf{Predictores físicos (5):} Frecuencia, ángulo de ataque, longitud de cuerda, velocidad de flujo libre y espesor de desplazamiento.
\item \textbf{Variable objetivo:} Nivel de presión sonora (SPL) en decibeles (dB).
\vspace{0.3cm}
\item \textbf{Hallazgos clave:}
\begin{itemize}
\item Relación claramente no lineal (forma de campana) entre frecuencia y SPL.
\item Alta multicolinealidad detectada en el espesor mediante el Factor de Inflación de Varianza (VIF), justificando el uso de regularización.
\end{itemize}
\end{itemize}
\end{frame}

\begin{frame}{Preprocesamiento y Pipeline}
\begin{block}{Transformaciones Físicamente Motivadas}
Se aplicó una transformación logarítmica ($\log_{10}$) a las variables de magnitud amplia (frecuencia, cuerda, velocidad y espesor) para capturar las tendencias lineales subyacentes asociadas a la acústica.
\end{block}

\vspace{0.3cm}

\begin{block}{Estandarización Aislada}
    Se integró un \textit{StandardScaler} dentro del proceso. Sus parámetros estadísticos ($\mu$, $\sigma$) se calcularon \textbf{estrictamente sobre el conjunto de entrenamiento} en cada pliegue para evitar que información futura contaminara el modelo.
\end{block}


\end{frame}

\begin{frame}{Esquema de Particionamiento (Prevención de Fuga de Información)}
\textbf{Problema:} La validación cruzada aleatoria simple dividiría la misma configuración de túnel entre entrenamiento y prueba, provocando resultados inflados. \~\

\textbf{Solución: Particiones Agrupadas}
\begin{itemize}
    \item Se construyó un identificador agrupando configuraciones únicas a partir de la triada \textit{cuerda + velocidad + ángulo}.
    \item Se reemplazó el \textit{KFold} estándar por \textbf{GroupKFold}.
    \item \textbf{Impacto:} El modelo se ve forzado a realizar estimaciones sobre configuraciones aerodinámicas que jamás observó durante su entrenamiento, reportando un error acústico realista.
\end{itemize}


\end{frame}

\end{document}